\documentstyle[%


righttag%


,11pt%


,amssymb%


,russian%


,izvru%


]{amsart}





%       Math definitions





\newcommand{\A}{{\cal A}}


\newcommand{\col}{\operatorname{col}}


\newcommand{\vraisup}{\operatorname{vrai~sup}}





\setcounter{page}{69}


\god{1997}


\nomer{5~(420)}


\udk{517.95}


\author{\it В.А.\,Севастьянов}


\title{МЕТОД РИМАНА ДЛЯ ТРЕХМЕРНОГО ГИПЕРБОЛИЧЕСКОГО\\


УРАВНЕНИЯ ТРЕТЬЕГО ПОРЯДКА}





%//\thanks{Работа выполнена при поддержке РФФИ, грант \No. 94-01-00183.}


\date{}


\pagestyle{plain}


\begin{document}





\maketitle





% Body text








{\bf 1.} Хорошо известен (\cite{Sob}, гл.\,5, \S\,3, c.\,68; \cite{Tri},


гл.\,3, \S\,3.4, с.\,195; \cite{Bi}, гл.\,1,  \S\,4, с.\,67 и др.)


классический метод Римана получения формулы явного решения задачи Коши


для уравнения


$$


u_{xy}+au_{x}+bu_{y}+cu=f,


$$


устанавливающий одновременно существование, единственность и непрерывную


зависимость этого решения от граничных условий. Здесь предлагается


ана\-ло\-гич\-ный метод для уравнения


\begin{equation}


L(u)\equiv u_{xyz}+au_{xy}+bu_{yz}+cu_{xz}+du_{x}+eu_{y}+fu_{z}+gu=F.


\end{equation}





{\bf 2.} В ориентированном системой координат $(x,y,z)$ пространстве


$\Bbb R^3$ рассмотрим поверхность $S$ класса


$C^{3}$, заданную уравнениями


$$


\begin{cases}


x = x(p,q), &\\


y = y(p,q), &\\


z = z(p,q), &(p,q)\in G^{2}\subset R^{2},


\end{cases}


\qquad \text{rang}\,


\begin{pmatrix}


x_{p} & y_{p} & z_{p} \\


x_{q} & y_{q} & z_{q}


\end{pmatrix}


=2.


$$





Предполагаем, что $S$ в каждой своей точке имеет касательную плоскость, не


параллельную ни одной из координатных осей, например,


$z_x>0$, $z_y>0$ в $G$. Проведем через точку $M(x_{0},y_{0},z_{0})$


плоскости $x=x_{0}$, $y=y_{0}$, $z=z_{0}$,


пересекающие $S$ по кривым $QC$, $CP$ и $PQ$ соответственно. Обозначим


$\Omega $ конечную область, ограниченную этими плоскостями и участком


$QCP$ поверхности $S$. Считаем ориентацию $\Omega $ положительной.





{\bf Постановка задачи:} {\it найти решение уравнения {\em (1)} класса


$C^{3}(\overline{\Omega } )$, удовлетворяющее граничным условиям}


\begin{equation}


\left.\frac{\partial^{k} u}{\partial l^{k}} \right\vert_{S} = \psi_{k},


\qquad k=0,1,2,


\end{equation}


{\it где $l$ --- заданное на $S$ некасательное к этой


поверхности поле направлений.}





Тем самым будет получено регулярное решение задачи (1)--(2) во всей области


$\Omega$.





{\bf 3.} Зададим поле направлений $l$ вектором $\Vec{a}


(l_{1}(p,q),l_{2}(p,q),l_{3}(p,q))$, $\Vec{a}\in C^{3}(G^{2})$,


причем $|\Vec{a}|\equiv 1$. Введем систему координат, связанную


с поверхностью $S$


\begin{equation}


\begin{cases}


x=x(p,q)+l_{1}(p,q)l, \\


y=y(p,q)+l_{2}(p,q)l, \\


z=z(p,q)+l_{3}(p,q)l, \quad l\in R.


\end{cases}


\end{equation}





В некоторой окрестности поверхности $S$ (3) определяет допустимое


преобразование координат, т.\,к. поле направлений $l$ по условию не


касательно к $S$ (\cite{So}, с.65). Следовательно, существует обратное


преобразование


$$


\begin{cases}


p=p(x,y,z),\\


q=q(x,y,z),\\


l=l(x,y,z)


\end{cases}


$$


класса $C^{3}$ в окрестности поверхности $S$ (\cite{Zo1}, с.488).





Определим функцию Римана $R(x,y,z,x_{0},y_{0},z_{0})$ уравнения (1) как


решение класса $C^{3}(\overline{\Omega} \times \overline{\Omega} )$


интегрального уравнения


\begin{gather}


v(x,y,z)-\int_{z_{0}}^{z} a(x,y,\zeta )v(x,y,\zeta )\, d\zeta


-\int_{x_{0}}^{x} b(\xi ,y,z)v(\xi ,y,z)\, d\xi -\notag\\


%


-\int_{y_{0}}^{y} c(x,\eta ,z)v(x,\eta,z)\, d\eta +


+\int_{y_{0}}^{y} \int_{z_{0}}^{z} d(x,\eta ,\zeta )v(x,\eta ,\zeta)\, d\zeta \, d\eta+\notag\\


%


+\int_{x_{0}}^{x} \int_{z_{0}}^{z} e(\xi ,y,\zeta )v(\xi ,y,\zeta )\,d\zeta \, d\xi +


+\int_{x_{0}}^{x} \int_{y_{0}}^{y} f(\xi ,\eta ,z)v(\xi ,\eta ,z)\, d\eta \, d\xi-\notag\\


%


-\int_{x_{0}}^{x} \int_{y_{0}}^{y} \int_{z_{0}}^{z} g(\xi ,\eta ,\zeta )v(\xi ,\eta


,\zeta )\, d\zeta \, d\eta \, d\xi =1,


\end{gather}


которое существует (\cite{Mu}, с.180).





\begin{thm} Пусть $\psi_{k} \in C^{3-k}(\overline{S} );$


$a,b,c\in C^{2}(\overline{\Omega} );$ $d,e,f\in C^{1}(\overline{\Omega} );$


$\ g,F\in C(\overline{\Omega} ).$ Тогда задача {\em (1)--(2)} в


$C^{3}(\overline{\Omega } )$


поставлена корректно, а ее решение имеет вид


\begin{gather}


u(x_{0},y_{0},z_{0})=\frac{(uR)_{C}+(uR)_{Q}+(uR)_{P}}{3} -\frac{1}{6} \biggl( \int_{QC} \Bigl \{


u_{z}R+u(3aR-2R_{z})\Bigr \} \, dz-\notag\\


%


\{ u_{y}R+u(3cR-2R_{y})\}dy+\int_{CP} \Bigl \{ u_{x}R+u(3bR-2R_{x})\Bigr \} \, dx-\notag\\


%


-\{u_{z}R+u(3aR-2R_{z})\}dz+\int_{PQ} \Bigl \{


u_{y}R+u(3cR-2R_{y})\Bigr \} \, dy-\{ u_{x}R+u(3bR-2R_{x})\} \, dx+\notag\\


%


+\iint_{PQC} \Bigl \{2u_{yz}R+


u_{y}(3aR-R_{z})+u_{z}(3cR-R_{y})+u(2R_{yz}-3(aR)_{y}-3(cR)_{z}+6dR)\Bigr \} \, dy\, dz+\notag\\


%


+\Bigl \{


2u_{xz}R+


+u_{x}(3aR-R_{z})+u_{z}(3bR-R_{x})+u(2R_{xz}-3(aR)_{x}-3(bR)_{z}+6eR)\Bigr \} \, dx\, dz+\notag\\


%


\Bigl \{2u_{xy}R+


u_{x}(3cR-R_{y})+u_{y}(3bR-R_{x})+u(2R_{xy}-3(bR)_{y}-3(cR)_{x}+6fR)\Bigr \}\,


dx\,dy\biggr) -\notag\\


%


-\int_{x_{0}}^{x_{P}} dx\int_{y_{0}}^{y_{PQ}(x)}dy\int_{z_{0}}^{z(x,y)} RF\, dz,


\end{gather}


где


\begin{gather}


u=\psi_{0},\quad u_{x}=p_{x}\psi_{0p}+q_{x}\psi_{0q}+l_{x}\psi_{1},\notag\\


%


u_{y}=p_{y}\psi_{0p}+q_{y}\psi_{0q}+l_{y}\psi_{1},\quad


u_{z}=p_{z}\psi_{0p}+q_{z}\psi_{0q}+l_{z}\psi_{1},\notag\\


%


u_{xy}=p_{x}p_{y}\psi_{0pp}+q_{x}q_{y}\psi_{0qq}+l_{x}l_{y}\psi_{2}


+(p_{x}q_{y}+q_{x}p_{y})\psi_{0pq}+\notag\\


%


+(p_{x}l_{y}+ l_{x}p_{y})\psi_{1p}+


(q_{x}l_{y}+l_{x}q_{y})\psi_{1q}+p_{xy}\psi_{0p}


+q_{xy}\psi_{0q}+l_{xy}\psi_{1},\notag\\


%


u_{xz}=p_{x}p_{z}\psi_{0pp}+q_{x}q_{z}\psi_{0qq}+l_{x}l_{z}\psi_{2}


+(p_{x}q_{z}+q_{x}p_{z})\psi_{0pq}+\notag\\


%


+(p_{x}l_{z}+ l_{x}p_{z})\psi_{1p}+


(q_{x}l_{z}+l_{x}q_{z})\psi_{1q}+p_{xz}\psi_{0p}


+q_{xz}\psi_{0q}+l_{xz}\psi_{1},\notag\\


%


u_{yz}=p_{y}p_{z}\psi_{0pp}+q_{y}q_{z}\psi_{0qq}+l_{y}l_{z}\psi_{2}


+(p_{y}q_{z}+q_{y}p_{z})\psi_{0pq}+\notag\\


%


+(p_{y}l_{z}+ l_{y}p_{z})\psi_{1p}+


(q_{y}l_{z}+l_{y}q_{z})\psi_{1q}+p_{yz}\psi_{0p}


+q_{yz}\psi_{0q}+l_{yz}\psi_{1}.


\end{gather}


\end{thm}





\begin{pf} Для начала отметим существование верхних пределов


интегрирования последнего интеграла в (5), т.\,к. все определители


$$


\begin{vmatrix}


x_{p} & y_{p} \\


x_{q} & y_{q}


\end{vmatrix}


\qquad


\begin{vmatrix}


x_{p} & z_{p} \\


x_{q} & z_{q}


\end{vmatrix}


\qquad


\begin{vmatrix}


y_{p} & z_{p} \\


y_{q} & z_{q}


\end{vmatrix}


$$


отличны от нуля в силу свойств поверхности $S$.





Для любой функции $u$ класса $C^{3}$ рассмотрим тождество


\begin{gather*}


RL(u)=\frac{\partial }{\partial x} \biggl( \frac{1}{3} (uR)_{yz}-\frac{1}{2} \left[


u(R_{z}-aR)\right] _{y}-\frac{1}{2} \left[ u(R_{y}-cR)\right] _{z}+\\


%


+ u\left[


R_{yz}-(aR)_{y}-(cR)_{z}+dR\right] \biggr) + \frac{\partial }{\partial y}


\biggl( \frac{1}{3} (uR)_{xz}-\frac{1}{2} \left[ u(R_{z}-aR)\right] _{x}-\\


%


-\frac{1}{2} \left[ u(R_{x}-bR)\right] _{z}+u\left[


R_{xz}-(aR)_{x}-(bR)_{z}+eR\right] \biggr) +\frac{\partial }{\partial z} \biggl(


\frac{1}{3} (uR)_{xy}-\\


%


-\frac{1}{2} \left[ u(R_{x}-bR)\right] _{y}-\frac{1}{2}


\left[ u(R_{y}-cR)\right] _{x}+u\left[ R_{xy}-(bR)_{y}-(cR)_{x}+fR\right]


\biggr),


\end{gather*}


представляющее лишь иную запись тождества (3) из статьи \cite{Zhe}. Интегрируя это


тождество по области $\Omega$, считая в нем $u$ решением уравнения (1), и


используя формулу Гаусса-Остроградского (\cite{Zo2}, гл.\,13, \S\,3,


с.241), получим


\begin{gather}


\iiint_{\Omega} RF\, dx\, dy\, dz=\iint_{\partial \Omega}


\biggl( \frac{1}{3} (uR)_{yz}-\frac{1}{2} \left[ u(R_{z}-aR)\right]


_{y}-\frac{1}{2} \left[ u(R_{y}-cR)\right]_{z}+\notag \\


%


+u\left[ R_{yz}-(aR)_{y}-(cR)_{z}+dR\right] \biggr)\, dy\wedge dz+


\biggl( \frac{1}{3} (uR)_{xz}-\frac{1}{2} \left[ u(R_{z}-aR)\right] _{x}-\notag\\


%


-\frac{1}{2} \left[ u(R_{x}-bR)\right] _{z}+u\left[


R_{xz}-(aR)_{x}-(bR)_{z}+eR\right] \biggr) \, dz\wedge dx+\notag\\


%


\biggl( \frac{1}{3}


(uR)_{xy}-


\frac{1}{2} \left[ u(R_{x}-bR)\right] _{y}-\frac{1}{2} \left[ u(R_{y}-cR)\right]_{x}+\notag\\


%


u\left[ R_{xy}-(bR)_{y}-(cR)_{x}+fR\right] \biggr) \, dx\wedge dy.


\end{gather}


Здесь знак ``$\wedge$'' --- внешнее умножение дифференциальных форм. Заменяя в


правой части, которую обозначим через $I$, интеграл по границе


$\partial\Omega$ суммой интегралов по ее составляющим $CQM$,


$PCM,$ $PQM$ и $PQC$ и учитывая тождества


\begin{multline*}


R_{yz}(x_{0},y,z,x_{0},y_{0},z_{0})-\bigl(


a(x_{0},y,z)R(x_{0},y,z,x_{0},y_{0},z_{0})\bigr)_{y}-\\


%


-\bigl( c(x_{0},y,z)R(x_{0},y,z,x_{0},y_{0},z_{0})\bigr)_{z}+


d(x_{0},y,z)R(x_{0},y,z,x_{0},y_{0},z_{0})\equiv 0,


\end{multline*}


\begin{multline*}


R_{xz}(x,y_{0},z,x_{0},y_{0},z_{0})-\bigl(


a(x,y_{0},z)R(x,y_{0},z,x_{0},y_{0},z_{0})\bigr)_{x}-\\


%


-\bigl( b(x,y_{0},z)R(x,y_{0},z,x_{0},y_{0},z_{0})\bigr)_{z}+


e(x,y_{0},z)R(x,y_{0},z,x_{0},y_{0},z_{0})\equiv 0,


\end{multline*}


\begin{multline*}


R_{xy}(x,y,z_{0},x_{0},y_{0},z_{0})-\bigl(


b(x,y,z_{0})R(x,y,z_{0},x_{0},y_{0},z_{0})\bigr)_{y}-\\


-\bigl( c(x,y,z_{0})R(x,y,z_{0},x_{0},y_{0},z_{0})\bigr)_{x}+


f(x,y,z_{0})R(x,y,z_{0},x_{0},y_{0},z_{0})\equiv 0,


\end{multline*}


получаемые подстановкой в уравнение (4) по очереди $x=x_{0},$


$y=y_{0},$ $z=z_{0}$


и дифференцированием его по $yz,$ $xz$ и $xy$ соответственно, имеем


\begin{gather*}


\displaybreak[3]


I=\iint_{CQM} \biggl\{ \frac{\partial }{\partial z} \biggl( \frac{1}{6}


(uR)_{y}-\frac{1}{2} u(R_{y}-cR) \biggr) +\frac{\partial }{\partial y} \biggl(


\frac{1}{6} (uR)_{z}-\frac{1}{2} u(R_{z}-aR)\biggr) \biggr\} \, dy\wedge dz+\\


%


+\iint_{PCM} \biggl\{ \frac{\partial }{\partial x} \biggl ( \frac{1}{6}


(uR)_{z}-\frac{1}{2} u(R_{z}-aR) \biggr ) +\frac{\partial }{\partial z} \biggl (


\frac{1}{6} (uR)_{x}-\frac{1}{2} u(R_{x}-bR) \biggr ) \biggr \} \, dz\wedge dx+\\


%//


%


+\iint_{PQM} \biggl \{ \frac{\partial }{\partial y} \biggl ( \frac{1}{6}


(uR)_{x}-\frac{1}{2} u(R_{x}-bR) \biggr ) +\frac{\partial }{\partial x} \biggl (


\frac{1}{6} (uR)_{y}-\frac{1}{2} u(R_{y}-cR)\biggr ) \biggr \} \, dx\wedge dy+\\


%


+\iint_{PQC} \biggl ( \frac{1}{3} u_{yz}R-\frac{1}{6} u_{y}R_{z}-\frac{1}{6} u_{z}R_{y}+\frac{1}{3} uR_{yz}+\frac{1}{2}


u_{y}aR-\frac{1}{2} u(aR)_{y}+\\


%


+\frac{1}{2} u_{z}cR-\frac{1}{2}u(cR)_{z}+u\,dR\biggr) \, dy\wedge dz+


\biggl( \frac{1}{3} u_{xz}R-\frac{1}{6} u_{x}R_{z}-\frac{1}{6} u_{z}R_{x}+\frac{1}{3} uR_{xz}+\\


+


+\frac{1}{2} u_{x}aR-\frac{1}{2} u(aR)_{x}+\frac{1}{2} u_{z}bR-\frac{1}{2}


u(bR)_{z}+ueR\biggr ) \, dz\wedge dx+\\


%


\biggl ( \frac{1}{3} u_{xy}R-\frac{1}{6}


u_{x}R_{y}-\frac{1}{6} u_{y}R_{x}+\frac{1}{3} uR_{xy}+\frac{1}{2}


u_{y}bR-\frac{1}{2} u(bR)_{y}+


\frac{1}{2} u_{x}cR-\frac{1}{2} u(cR)_{x}+ufR\biggr) \, dx\wedge dy.


\end{gather*}





По формуле Грина (\cite{Zo2}, гл.\,13, \S\,3, с.236) все интегралы по


плоским областям сводятся к однократным


интегралам по замкнутым контурам, причем любой участок каждого контура


обходится дважды --- в противоположных направлениях, как это и должно быть


у ориентированной поверхности:


\begin{gather*}


I =\int_{CQM} \biggl( \frac{1}{6} (uR)_{z}-\frac{1}{2} u(R_{z}-aR)\biggr) \,


dz-\biggl( \frac{1}{6} (uR)_{y}-\frac{1}{2} u(R_{y}-cR)\biggr) \, dy+\\


%


 +\int_{PCM} \biggl( \frac{1}{6} (uR)_{x}-\frac{1}{2} u(R_{x}-bR)\biggr) \,


dx-\biggl( \frac{1}{6} (uR)_{z}-\frac{1}{2} u(R_{z}-aR)\biggr) \, dz+\\


%


 +\int_{PQM} \biggl( \frac{1}{6} (uR)_{y}-\frac{1}{2} u(R_{y}-cR)\biggr) \,


dy-\biggl( \frac{1}{6} (uR)_{x}-\frac{1}{2} u(R_{x}-bR)\biggr) \, dx+\\


%


 +\frac{1}{6} \iint_{PQC} \biggl \{


2u_{yz}R+u_{y}(3aR-R_{z})+u_{z}(3cR-R_{y})+u(2R_{yz}-3(aR)_{y}-3(cR)_{z}+


6dR)\biggr \} \, dy\wedge dz+\\


 +\Bigl \{ 2u_{xz}R+u_{x}(3aR-R_{z})+u_{z}(3bR-


R_{x})+u(2R_{xz}-3(aR)_{x}-


3(bR)_{z}+6eR)\Bigr \} \, dz\wedge dx+\\


 +\Bigl \{


2u_{xy}R+u_{x}(3cR-\Bigr. R_{y})+u_{y}(3bR-R_{x})+u(2R_{xy}-


3(bR)_{y}-3(cR)_{x}+6fR)\Bigr \} \,dx\wedge dy.


\end{gather*}





Аналогично случаю с 2-кратными интегралами,


заменяя каждый криволинейный интеграл суммой интегралов по составляющим


его контура и используя тождества


\begin{align*}


R_{z}(x_{0},y_{0},z,x_{0},y_{0},z_{0})- a(x_{0},y_{0},z)R(x_{0},y_{0},z,x_{0},y_{0},z_{0})&\equiv 0,\\


R_{y}(x_{0},y,z_{0},x_{0},y_{0},z_{0})- c(x_{0},y,z_{0})R(x_{0},y,z_{0},x_{0},y_{0},z_{0})&\equiv 0,\\


R_{x}(x,y_{0},z_{0},x_{0},y_{0},z_{0})- b(x,y_{0},z_{0})R(x,y_{0},z_{0},x_{0},y_{0},z_{0})&\equiv 0,


\end{align*}


непосредственно вытекающие из определения функции Римана, получим, что


сумма этих однократных интегралов перепишется:


\begin{gather*}


K=\frac{1}{6} \int_{CM} (uR)_{z}\, dz-\frac{1}{6} \int_{MQ}


(uR)_{y}\, dy+\frac{1}{6} \int_{QC} (u_{z}R+uR_{z}-3uR_{z}+3uaR)\, dz-\\


%


-(u_{y}R+uR_{y}-3uR_{y}+3ucR)\, dy+\frac{1}{6} \int_{PM} (uR)_{x}\,dx


-\frac{1}{6} \int_{MC} (uR)_{z}\, dz+\\


%


+\frac{1}{6} \int_{CP}(u_{x}R+uR_{x}-


3uR_{x}+3ubR)\, dx-(u_{z}R+uR_{z}-3uR_{z}+3uaR)\,dz+\\


%


+\frac{1}{6} \int_{QM} (uR)_{y}\, dy-\frac{1}{6} \int_{MP}


(uR)_{x}\, dx+\frac{1}{6} \int_{PQ} (u_{y}R+uR_{y}-3uR_{y}+3ucR)\, dy-\\


-(u_{x}R+uR_{x}-3uR_{x}+3ubR)\, dx.


\end{gather*}





Вычисляя здесь  интегралы по отрезкам прямых и используя представления


$K$ и $I$, из (7) окончательно получаем формулу (5). Представления


производных решения $u$ на $S$ (6) получаются дифференцированием


$u(p(x,y,z)$, $q(x,y,z),$ $l(x,y,z))$ как сложной функции.





Формула (5) получена в предположении существования решения исходной


задачи, т.\,е. если решение существует, то оно единственно. Непосредственной


под\-ста\-нов\-кой выражения (5) в уравнение (1) можно  убедиться,  что  оно


является решением этого уравнения. При этом используются формулы


дифференцирования интегралов, зависящих от параметра, свойство функции


Римана $R$ и ее производных $R_{x},R_{xy},\dotsc$ быть решением уравнения


$L(u)=0$ по второй тройке аргументов. Доказательство этого свойства аналогично


2-мерному случаю (\cite{Vek}, с.25) на основе представления решения 3-х мерной


задачи Гурса \cite{Zhe}. Тем самым обеспечено существование решения поставленной


задачи.





Непрерывная зависимость решения $u$ от $\psi_{k}$ в $C^3(\overline{\Omega} )$,


$k=0,1,2$, следует непосредственно из (5) в результате элементарных


оценок в терминах $\varepsilon$, $\delta$.


\end{pf}





{\bf 4.} В случаях, когда гладкость начальных данных нарушается, трижды


непрерывно дифференцируемых решений задачи, вообще говоря, не существует,


и решение уравнения (1) следует понимать в обобщенном смысле.





Если начальные данные задачи Коши (1)--(2) имеют, например, в точке


$P$ особенность, то решение задачи Коши будет иметь особенности вдоль


прямых, проходящих через точку $P$ и параллельных координатным осям.


То же справедливо для точек $Q$ и $C$.





Уравнение (1) является обобщением 2-мерного уравнения Лапласа, имеющего


приложения в теории поверхностей (\cite{Mu}, с.158).





\begin{thebibliography}{99}





\bibitem{Sob}


Соболев С.Л. {\em Уравнения математической физики.} -- 5-е изд. -- М.:


Наука, 1992. -- 432 с.





\bibitem{Tri}


Трикоми Ф. {\em Лекции по уравнениям в частных производных.} -- М.: Ин. лит.,


1957. -- 443 с.





\bibitem{Bi}


Бицадзе А.В. {\em Некоторые классы уравнений в частных производных.}


-- М.: Наука, 1981. -- 448~с.





\bibitem{So}


Сокольников И.С.  {\em Тензорный анализ. Теория и применения в геометрии


и в механике сплошных сред.} -- М.: Наука, 1971. -- 376 с.





\bibitem{Zo1}


Зорич В.А. {\em Математический анализ.} Ч.\,1. -- М.: Наука, 1981. -- 543 с.





\bibitem{Mu}


Мюнтц Г.М. {\em Интегральные уравнения.} Ч.1. -- Л.--М.: Гостехиздат,


1934. -- 320 с.





\bibitem{Zhe}


Жегалов В.И. {\em Трехмерный аналог задачи Гурса} // Неклассич. уравнения и


уравнения смеш. типа. -- Новосибирск, 1990. -- С.\,94--98.





\bibitem{Zo2}


Зорич В.А. {\em Математический анализ.} Ч.2. -- М.: Наука, 1984. -- 640 с.





\bibitem{Vek}


Векуа И.Н. {\em Новые методы решения эллиптических


уравнений.} -- М.--Л.: Гостехиздат, 1948. -- 296 с.





\end{thebibliography}





\vskip 2cm





\postup{\it Казанский государственный университет}{\it Поступили}


\postup{}{{\it первый вариант} 17.07.1995}


\postup{}{{\it окончательный вариант} 01.04.1996}





\end{document}


